% R2 — within-sepsis mortality prediction is null at late sampling windows (draft 2026-09-07)
\subsection{Within-sepsis, no model predicts mortality above chance at late sampling windows}
\label{ssec:r2_within_sepsis}

We next ran the identical protocol on the sepsis arm (91 patients, 17 fatal,
28-day mortality; Section~\ref{ssec:methods_cohort}). Sepsis is
life-threatening organ dysfunction caused by a dysregulated host response to
infection \citep{singer2016sepsis3} and accounts for roughly 11 million
deaths annually worldwide \citep{rudd2020global}. This is, to our
knowledge, the largest platelet single-cell sepsis cohort with per-patient
outcome labels assembled to date, and the first evaluated under patient-level
cross-validation. No model exceeded chance
(Table~\ref{tab:within_sepsis} and Figure~\ref{fig:within_sepsis}).
\model{UCE} scored highest (patient-level AUC 0.548 $\pm$ 0.038), followed by
\model{scPRINT-2} (0.523 $\pm$ 0.049), \model{Geneformer} (0.515 $\pm$ 0.043),
and \model{STATE} (0.511 $\pm$ 0.040), with \model{scGPT} (0.449 $\pm$ 0.052)
and the PCA-200 baseline (0.428 $\pm$ 0.082) both below chance. Under
patient-level paired bootstrap with multi-comparison control, no pairwise
difference between models survived:
the only nominal contrast was \model{UCE} over \model{scGPT} ($+0.101$,
95\% CI [0.001, 0.197], $p = 0.047$, not surviving correction), and no model
separated from the PCA-200 baseline (all 95\% confidence intervals include
zero; smallest nominal $p = 0.079$ for \model{UCE}). In other words, the
within-sepsis comparison is a complete null: not only does every method sit at
chance, no ordering among them can be defended at this sample size. At 17
fatal events the benchmark resolves only large effects: 80\% power to detect
a single predictor requires a patient-level AUC of about 0.71 or higher, so
this null bounds rather than excludes smaller signal.

The cohort's sampling times complicate any stronger reading
(Section~\ref{ssec:limitations}). Draws occurred a median of 16 days after
admission (range 3--121), and fatal cases were drawn earlier than survivors
(median 10 versus 18.5 days), so sampling time is associated with the outcome
and could, in principle, masquerade as signal. The measured result is a null
despite that confound working in the models' favor. The literature context
sharpens the interpretation: outcome-predictive platelet transcriptomic signals
in sepsis have so far been documented only in early windows, within hours of
diagnosis or within 72 hours of ICU admission
\citep{qiu2021sepsis, middleton2019sepsis, middleton2020nets}, against a
background of sepsis host-response heterogeneity established by single-cell
and transcriptomic studies \citep{reyes2020immune, davenport2016genomic,
sweeney2018unsupervised}. Our measurement is the first at
draws beyond the first days of illness, and it bounds when that signal remains
predictable: by the second week after admission, neither foundation-model
embeddings nor simple baselines carried usable patient-level mortality signal.
Because the early-window evidence comes from different cohorts and platforms,
this comparison is between studies; a within-cohort temporal test would
require serial draws we do not have.
